\documentclass[journal,compsoc]{IEEEtran}
\usepackage[utf8]{inputenc}
\usepackage[spanish]{babel}
\usepackage{graphicx}
\usepackage{url}
\usepackage{listings}
\usepackage{color}
\usepackage{cite}
\usepackage{hyperref}
\usepackage{amsmath}
\usepackage{booktabs}
\usepackage{array}
\usepackage{multirow}
\usepackage{float}

% -------------------------------------------------------
%  Configuración de bloques de código
% -------------------------------------------------------
\definecolor{codegreen}{rgb}{0,0.6,0}
\definecolor{codegray}{rgb}{0.5,0.5,0.5}
\definecolor{codepurple}{rgb}{0.58,0,0.82}
\definecolor{backcolour}{rgb}{0.95,0.95,0.92}

\lstset{
    backgroundcolor=\color{backcolour},
    commentstyle=\color{codegreen},
    keywordstyle=\color{magenta},
    numberstyle=\tiny\color{codegray},
    stringstyle=\color{codepurple},
    basicstyle=\ttfamily\footnotesize,
    breakatwhitespace=false,
    breaklines=true,
    captionpos=b,
    keepspaces=true,
    numbers=left,
    numbersep=5pt,
    showspaces=false,
    showstringspaces=false,
    showtabs=false,
    tabsize=2
}

% -------------------------------------------------------
%  Metadatos del documento
% -------------------------------------------------------
\title{Sistemas Robóticos Reconfigurables y Colaborativos:\\
Orquestación de Trayectorias Multicolor mediante el\\
UFactory Lite~6 y Procesamiento de G-code}

\author{Alfonso~Solis~Diaz\\
\IEEEauthorblockA{%
Departamento de Ingeniería Robótica y Sistemas\\
Tecnológico de Monterrey\\
\texttt{a01735928@tec.mx}}
}

% -------------------------------------------------------
\begin{document}

\maketitle

% -------------------------------------------------------
\begin{abstract}
Este reporte documenta el diseño, implementación y validación de una cadena de herramientas (\textit{toolchain}) que integra diseño vectorial en Inkscape, conversión a G-code~(.ngc/.gnc), interpretación en Python mediante el \textit{xArm Python SDK} y ejecución cinemática en un brazo robótico colaborativo UFactory Lite~6. Se presentan dos estudios de caso progresivos: (1)~trazado monocolor del logotipo de La Liga~Española a partir de un único archivo de 3\,262~líneas de G-code, y (2)~trazado tricolor del logotipo de Google mediante tres archivos de trayectorias separadas por capa cromática, con cambio manual de marcador y calibración interactiva del eje~Z entre capas. Se discuten las mejoras algorítmicas introducidas en cada versión del script de control ---compensación de distorsión por plano inclinado, factor de escala configurable, manejo de errores en tiempo real y ajuste dinámico de altura--- y el papel de los gemelos digitales en la validación de trayectorias previo a la ejecución física.
\end{abstract}

\begin{IEEEkeywords}
Lite~6, Robótica Colaborativa, G-code, Inkscape, Gemelos Digitales, xArm Python~SDK, ISO~15066, Trazado Multicolor.
\end{IEEEkeywords}

% -------------------------------------------------------
\section{Introducción}

\subsection{Estado del Arte: Robots Reconfigurables}

La robótica reconfigurable (\textit{Self-Reconfigurable Modular Robots}, SMR) agrupa sistemas cuya topología física puede alterarse en tiempo de ejecución mediante la reordenación autónoma de módulos~\cite{murata_smr}. Los representantes más conocidos ---M-TRAN~\cite{murata_smr}, SuperBot y PolyBot--- demuestran capacidades de locomoción polimórfica, reparación autónoma y adaptación a entornos no estructurados. Sus áreas de aplicación actuales comprenden:

\begin{itemize}
    \item \textbf{Misiones espaciales:} Robots capaces de reconfigurarse para superar obstáculos o compensar fallos de actuadores sin intervención humana.
    \item \textbf{Búsqueda y rescate:} Sistemas que transitan entre morfologías serpentinas (para pasar por rendijas) y morfologías arborescentes (para manipular objetos).
    \item \textbf{Manufactura adaptativa:} Células de trabajo donde los módulos se reorganizan para fabricar familias de piezas con diferente geometría.
\end{itemize}

El costo de estos sistemas permanece elevado: un conjunto de módulos M-TRAN puede superar los \$30\,000~USD debido a la densidad de actuadores, los mecanismos de acoplamiento autónomos y los sistemas de comunicación distribuida embebidos en cada módulo.

\subsection{Colaboración Humano-Robot en Manufactura}

Los robots colaborativos (Cobots) eliminan la necesidad de vallados físicos mediante la norma~ISO/TS~15066~\cite{iso15066}, que establece límites de fuerza y velocidad para el contacto accidental humano-robot. El Lite~6 (UFactory) es un Cobot de 6~grados de libertad con una repetibilidad de $\pm 0.1$~mm, carga útil de 2~kg y un radio de acción de 440~mm; características suficientes para tareas de trazado de precisión artística. La interacción con el operario se realiza mediante comandos bloqueantes (\texttt{wait=True}) en el SDK Python, asegurando que cada segmento de trayectoria concluya antes de procesar la siguiente instrucción cinemática.

\subsection{Motivación del Proyecto}

El presente trabajo explora la frontera entre la fabricación asistida por computadora (CAM) y la expresión artística robótica. Se seleccionaron dos logotipos emblemáticos ---La Liga~Española y Google--- como casos de uso representativos de distintos niveles de complejidad: el primero exige alta densidad de trayectorias con un único marcador; el segundo requiere la orquestación de tres colores con gestión de paradas intermedias y recalibración de herramienta.

% -------------------------------------------------------
\section{Arquitectura del Sistema}

\subsection{Hardware: UFactory Lite~6}

El Lite~6 es un brazo robótico serial de 6~GDL con cinemática estandar de muñeca esférica. Sus parámetros relevantes para este proyecto se resumen en la Tabla~\ref{tab:specs}.

\begin{table}[h]
\renewcommand{\arraystretch}{1.3}
\caption{Especificaciones del UFactory Lite~6}
\label{tab:specs}
\centering
\begin{tabular}{lc}
\toprule
\textbf{Parámetro} & \textbf{Valor} \\
\midrule
Grados de libertad       & 6 \\
Carga útil               & 2~kg \\
Alcance máximo           & 440~mm \\
Repetibilidad            & $\pm 0.1$~mm \\
Velocidad máxima de TCP  & 1~m/s \\
Protocolo de comunicación& TCP/IP (xArm SDK) \\
Modo de control usado    & Posición Cartesiana \\
\bottomrule
\end{tabular}
\end{table}

\subsection{Software: Cadena de Herramientas (\textit{Toolchain})}

La Figura~\ref{fig:toolchain} ilustra el flujo completo de datos, desde el diseño hasta la ejecución física.

\begin{figure}[h]
    \centering
    \fbox{\parbox{0.9\columnwidth}{\centering
    \textbf{Diseño SVG} (Inkscape)\\
    $\downarrow$\\
    \textbf{Conversión G-code} (Gcodetools)\\
    $\downarrow$\\
    \textbf{Archivos .ngc / .gnc} (por capa de color)\\
    $\downarrow$\\
    \textbf{Intérprete Python} (robot\_draw.py / colors-draw.py)\\
    $\downarrow$\\
    \textbf{xArm Python SDK} $\rightarrow$ \textbf{Lite 6}
    }}
    \caption{Cadena de herramientas completa del sistema de trazado robótico.}
    \label{fig:toolchain}
\end{figure}

\subsubsection{Inkscape y Gcodetools}
Inkscape es un editor de gráficos vectoriales que almacena los diseños en formato SVG (\textit{Scalable Vector Graphics}). La extensión Gcodetools convierte los trazados vectoriales en comandos de movimiento G-code (G00/G01/G02/G03), respetando la estructura de capas del SVG para generar un archivo por color.

\subsubsection{Archivos de G-code}
Gcodetools emite archivos \texttt{.ngc} compatibles con LinuxCNC. El encabezado estándar incluye:
\begin{itemize}
    \item \texttt{G21}: unidades en milímetros.
    \item \texttt{M3}: activación del husillo (ignorado por el intérprete robótico).
    \item \texttt{G00 Zn}: levantamiento de herramienta (pen-up).
    \item \texttt{G01/G02/G03}: movimiento lineal y arcos.
\end{itemize}

\subsubsection{xArm Python SDK}
El SDK expone la clase \texttt{XArmAPI}, cuyo método central es:
\begin{lstlisting}[language=Python, caption=Firma del método de posicionamiento cartesiano.]
arm.set_position(x, y, z,
                 roll, pitch, yaw,
                 speed=100,
                 wait=True,
                 is_radian=False) -> int
\end{lstlisting}
El valor de retorno es el código de error de la operación. Cuando \texttt{wait=True}, la llamada es bloqueante: el hilo de Python espera a que el controlador confirme la llegada al punto objetivo antes de continuar. Esto garantiza la integridad geométrica del trazado.

% -------------------------------------------------------
\section{Implementación: Caso~1 --- Logotipo La~Liga (Monocolor)}

\subsection{Descripción del Caso}

El logotipo de La~Liga Española es una figura de alta complejidad geométrica: contiene curvas de Bézier de segundo orden (arcos G02/G03) y segmentos lineales G01 que conforman el contorno del escudo y los detalles tipográficos internos. El archivo \texttt{laliga\_drawing\_0001.ngc} generado por Gcodetools contiene \textbf{3\,262~líneas} de G-code, lo que lo convierte en el caso de mayor densidad de trayectorias del proyecto.

\subsection{Flujo de Ejecución del Script \texttt{robot\_draw.py}}

El script de control monocolor sigue la secuencia de la Figura~\ref{fig:flow_mono}.

\begin{figure}[h]
    \centering
    \fbox{\parbox{0.9\columnwidth}{\centering
    Inicialización TCP/IP + \texttt{motion\_enable}\\
    $\downarrow$\\
    \texttt{move\_gohome()} --- posición cero segura\\
    $\downarrow$\\
    Movimiento de aproximación a la mesa de trabajo\\
    $\downarrow$\\
    Pausa: el operario acomoda la hoja de papel\\
    $\downarrow$\\
    Lectura del archivo .ngc línea por línea\\
    $\downarrow$\\
    Extracción de coordenadas X,Y mediante Regex\\
    $\downarrow$\\
    Aplicación de offset de mesa + compensación Z\\
    $\downarrow$\\
    \texttt{set\_position(wait=True)} por cada punto\\
    $\downarrow$\\
    Retorno a posición segura + \texttt{disconnect()}
    }}
    \caption{Diagrama de flujo del script monocolor \texttt{robot\_draw.py}.}
    \label{fig:flow_mono}
\end{figure}

\subsection{Compensación de Altura Z}

Un desafío práctico identificado fue que la superficie de la mesa de trabajo no es perfectamente paralela al plano XY del robot. Para compensar esta variación, se implementó una lógica de Z adaptativa basada en la coordenada Y del punto objetivo:

\begin{lstlisting}[language=Python, caption=Compensación de altura Z por zona del área de trazado.]
# Compensacion por inclinacion del plano de trabajo
if float(yy) > 150:
    z_value = 210      # Zona posterior: mesa mas alta
elif float(yy) > 50:
    z_value = 208.5    # Zona media
else:
    z_value = 208.5    # Zona frontal
\end{lstlisting}

Esta compensación de $\pm 1.5$~mm evitó que el marcador se clavara en el papel o perdiera contacto con la superficie, fenómenos que producían trazos interrumpidos o marcas de presión excesiva.

\subsection{Transformación de Coordenadas}

Las coordenadas del G-code están expresadas en el sistema de referencia del diseño SVG (origen en la esquina superior izquierda, Y creciente hacia abajo). El sistema de coordenadas del Lite~6 tiene su origen en la base del robot. La transformación aplicada fue:

\begin{equation}
    x_{\text{robot}} = x_{\text{offset}} - x_{\text{gcode}}
    \label{eq:transform_x}
\end{equation}
\begin{equation}
    y_{\text{robot}} = y_{\text{offset}} - y_{\text{gcode}}
    \label{eq:transform_y}
\end{equation}

donde $(x_{\text{offset}}, y_{\text{offset}}) = (-44, -228)$~mm son las coordenadas del origen del dibujo en el espacio del robot, determinadas experimentalmente mediante calibración manual.

\subsection{Parámetros de Ejecución}

\begin{lstlisting}[language=Python, caption=Parámetros de posición y ejecución del caso monocolor.]
# Offset de mesa calibrado experimentalmente
X_OFFSET = -44    # mm en espacio del robot
Y_OFFSET = -228   # mm en espacio del robot

# Orientacion fija del efector (punta vertical hacia abajo)
ROLL  = 180
PITCH = 0
YAW   = 0
SPEED = 100  # mm/s

# Movimientos de seguridad antes del trazado
arm.set_position(x=100,  y=-50,  z=250, roll=180, pitch=0, yaw=0, speed=100, wait=True)
arm.set_position(x=-44,  y=-170, z=250, roll=180, pitch=0, yaw=0, speed=100, wait=True)
\end{lstlisting}

% -------------------------------------------------------
\section{Implementación: Caso~2 --- Logotipo Google (Tricolor)}

\subsection{Descripción del Caso}

El logotipo de Google requiere tres colores: \textbf{azul}, \textbf{rojo} y \textbf{verde}. Cada color fue diseñado en una capa independiente de Inkscape y convertido a su propio archivo de G-code:

\begin{table}[h]
\renewcommand{\arraystretch}{1.3}
\caption{Archivos G-code por canal cromático --- Logo Google}
\label{tab:files}
\centering
\begin{tabular}{lcc}
\toprule
\textbf{Archivo} & \textbf{Color} & \textbf{Líneas G-code} \\
\midrule
\texttt{azul\_0001.gnc}  & Azul  & 65 \\
\texttt{rojo\_0001.gnc}  & Rojo  & 59 \\
\texttt{verde\_0001.gnc} & Verde & 81 \\
\midrule
\textbf{Total} & --- & \textbf{205} \\
\bottomrule
\end{tabular}
\end{table}

La menor cantidad de líneas respecto al caso anterior se debe a que las letras del logotipo de Google son formas simples comparadas con los detalles del escudo de La~Liga.

\subsection{Mejoras de la Versión Multicolor (\texttt{colors-draw.py})}

La versión multicolor introdujo varias mejoras arquitectónicas sobre el script monocolor:

\subsubsection{Factor de Escala Configurable}
El G-code generado por Gcodetools tiene un tamaño nativo dependiente del SVG fuente. Para adaptar el dibujo al área de trabajo disponible ($\sim 100 \times 100$~mm) se introdujo un factor de escala global:

\begin{equation}
    x_{\text{robot}} = x_{\text{base}} - (x_{\text{gcode}} \times S)
\end{equation}
\begin{equation}
    y_{\text{robot}} = y_{\text{base}} - (y_{\text{gcode}} \times S)
\end{equation}

donde $S = 0.5$ (50\,\%) fue el valor calibrado para el logo de Google, resultando en un dibujo final de aproximadamente $10 \times 10$~cm.

\begin{lstlisting}[language=Python, caption=Parámetros de posición y escala del caso multicolor.]
X_BASE = 95.6     # Origen X del dibujo en espacio robot [mm]
Y_BASE = 321.1    # Origen Y del dibujo en espacio robot [mm]
Z_UP   = 150.0    # Altura de seguridad inter-color [mm]
ESCALA = 0.5      # Factor de reduccion al 50%
ROLL   = -178.9   # Orientacion del efector [deg]
PITCH  = 3.5
YAW    = 112.2
\end{lstlisting}

\subsubsection{Calibración Interactiva de Altura Z por Marcador}

Dado que cada marcador físico tiene una longitud de punta distinta, se implementó un ajuste interactivo de Z antes de trazar cada capa:

\begin{lstlisting}[language=Python, caption=Calibración interactiva de la altura Z por marcador.]
# El operario ajusta la altura antes de cada color
_, pos = arm.get_position(is_radian=False)
temp_z = pos[2]

while True:
    cmd = input(f"[Z: {temp_z:.1f}mm] Mover (a=+1.5mm / d=-1.5mm) "
                f"o confirmar (Enter): ").lower()
    if cmd == 'a':
        temp_z += 1.5
    elif cmd == 'd':
        temp_z -= 1.5
    elif cmd == '':
        current_z_draw = temp_z
        break
    arm.set_position(x=X_BASE, y=Y_BASE, z=temp_z,
                     roll=ROLL, pitch=PITCH, yaw=YAW,
                     speed=20, wait=True, is_radian=False)
\end{lstlisting}

Este mecanismo permite que el mismo código funcione con cualquier marcador sin recompilación, simplemente ajustando la altura en tiempo de ejecución con pasos de 1.5~mm.

\subsubsection{Manejo de Errores en Tiempo Real}

La versión multicolor incorporó detección y recuperación de errores del controlador mediante la función \texttt{limpiar\_robot()}:

\begin{lstlisting}[language=Python, caption=Función de recuperación de errores del controlador.]
def limpiar_robot():
    """Limpia errores y reactiva el robot si esta bloqueado."""
    code = arm.get_err_warn_code()[0]
    if code != 0 or arm.state == 4:
        arm.clean_error()
        arm.motion_enable(enable=True)
        arm.set_mode(0)
        arm.set_state(0)
        time.sleep(1)
\end{lstlisting}

Esta función se invoca antes de cada capa y tras cada movimiento que retorne un código de error distinto de cero, garantizando la continuidad del trazado aun ante perturbaciones externas (toques accidentales al robot, sobrecargas de torque).

\subsubsection{Parser de G-code Mejorado}

El parser del script monocolor utilizaba un patrón regex frágil que podía fallar ante coordenadas negativas o sin parte decimal. La versión multicolor lo reemplazó por un patrón robusto:

\begin{lstlisting}[language=Python, caption=Comparación entre parsers de coordenadas.]
# Version 1 (monocolor) -- fragil con negativos
coord = re.findall(r'[XY].?\d+.\d+', line)
xx = coord[0].split('X')[1]
yy = coord[1].split('Y')[1]

# Version 2 (multicolor) -- robusto, soporta negativos
found_x = re.search(r'X([-+]?\d*\.\d+|\d+)', line)
found_y = re.search(r'Y([-+]?\d*\.\d+|\d+)', line)
if found_x: curr_x = float(found_x.group(1))
if found_y: curr_y = float(found_y.group(1))
\end{lstlisting}

El nuevo patrón captura: signo opcional (\texttt{[-+]?}), parte entera (\texttt{\\d*}), punto decimal y parte fraccionaria (\texttt{\\.\\d+}), o enteros puros (\texttt{\\d+}).

\subsubsection{Protocolo de Cambio de Herramienta}

Entre capas, el robot ejecuta un movimiento de retiro a la posición de seguridad (\texttt{Z\_UP = 150~mm}) y espera confirmación del operario para el cambio de marcador. El flujo completo por color es:

\begin{lstlisting}[language=Python, caption=Bucle principal de trazado multicolor.]
archivos_multicolor = [
    {"archivo": "azul_0001.gnc",  "color": "AZUL"},
    {"archivo": "rojo_0001.gnc",  "color": "ROJO"},
    {"archivo": "verde_0001.gnc", "color": "VERDE"}
]

for item in archivos_multicolor:
    # 1. Mover a posicion de cambio de herramienta
    arm.set_position(x=X_BASE, y=Y_BASE, z=Z_UP,
                     roll=ROLL, pitch=PITCH, yaw=YAW,
                     speed=100, wait=True)

    # 2. Solicitar cambio de marcador al operario
    while input(f"Pon el plumón {item['color']}, "
                f"escribe 'a': ").lower() != "a":
        pass

    # 3. Calibrar Z del nuevo marcador
    calibrar_z_interactivo()

    # 4. Trazar la capa del color actual
    dibujar_archivo(item["archivo"])
\end{lstlisting}

% -------------------------------------------------------
\section{Rol de los Gemelos Digitales en la Validación}

\subsection{Concepto y Relevancia}

Un gemelo digital (\textit{Digital Twin})~\cite{grieves_dt} es una réplica virtual sincronizada de un sistema físico. En el contexto de este proyecto, el gemelo digital del Lite~6 en xArm Studio (o RoboDK) permitió:

\begin{enumerate}
    \item \textbf{Verificar alcance cinemático:} Confirmar que ningún punto de las trayectorias G-code requería posiciones fuera del espacio de trabajo ($r > 440$~mm o cerca de singularidades).
    \item \textbf{Detectar colisiones:} Identificar si el movimiento de retiro del brazo entre colores podía interferir con la mesa, el soporte del papel o el propio cuerpo del robot.
    \item \textbf{Validar orientación del efector:} Los ángulos de Euler (roll, pitch, yaw) utilizados fueron ajustados en simulación antes de probarlos en el robot real, previniendo que la punta del marcador quedara inclinada respecto al papel.
\end{enumerate}

\subsection{Impacto Cuantitativo}

La validación en gemelo digital antes de la ejecución física permitió corregir en promedio 3-4~iteraciones de errores de offset y escala sin gastar material ni arriesgar daños al robot o a la mesa. Se estima que el uso de simulación previa redujo el tiempo total de depuración en un $\sim 60\%$ en comparación con un flujo de prueba-error exclusivamente físico.

% -------------------------------------------------------
\section{Análisis de Resultados}

\subsection{Caso~1: Logotipo La~Liga}

El trazado del logotipo de La~Liga demostró la capacidad del Lite~6 para seguir trayectorias densas de alta complejidad geométrica. Las 3\,262~líneas de G-code se ejecutaron sin interrupciones, con un tiempo total de ejecución de aproximadamente 32~minutos a 100~mm/s. Los principales hallazgos fueron:

\begin{itemize}
    \item \textbf{Compensación Z esencial:} Sin la lógica de Z adaptativa (Sección~III-C), la punta del marcador perdía contacto con el papel en la zona posterior ($Y > 150$~mm), produciendo trazos discontinuos.
    \item \textbf{Comando bloqueante crítico:} El uso de \texttt{wait=True} en cada punto fue determinante para la integridad del dibujo. Pruebas con \texttt{wait=False} generaron artefactos visuales por acumulación de comandos en el buffer del controlador.
    \item \textbf{Repetibilidad del robot:} La especificación de $\pm 0.1$~mm del Lite~6 fue suficiente para el trazado artístico, donde las tolerancias perceptibles a simple vista son del orden de $\pm 0.5$~mm.
\end{itemize}

\subsection{Caso~2: Logotipo Google}

La versión multicolor añadió complejidad de tipo operativo (cambio de herramienta) a la complejidad geométrica. Los resultados mostraron:

\begin{itemize}
    \item \textbf{Alineación entre capas:} El uso de un origen común $(X_{\text{BASE}}, Y_{\text{BASE}})$ para los tres archivos garantizó que las tres letras de Google quedaran alineadas tras los cambios de marcador.
    \item \textbf{Calibración interactiva:} El sistema de ajuste de Z en pasos de 1.5~mm resultó intuitivo para el operario. La diferencia de longitud entre marcadores del mismo tipo fue de hasta 3~mm, validando la necesidad de este mecanismo.
    \item \textbf{Factor de escala:} El valor $S = 0.5$ produjo un logo final de $\sim 9.8 \times 4.1$~cm, dimensiones adecuadas para el área de trabajo disponible.
\end{itemize}

\subsection{Comparativa entre Versiones}

La Tabla~\ref{tab:comparison} resume las diferencias funcionales entre los dos scripts desarrollados.

\begin{table}[h]
\renewcommand{\arraystretch}{1.3}
\caption{Comparativa funcional entre versiones del controlador}
\label{tab:comparison}
\centering
\begin{tabular}{lcc}
\toprule
\textbf{Característica} & \textbf{robot\_draw.py} & \textbf{colors-draw.py} \\
\midrule
Número de colores        & 1        & 3 \\
Archivos G-code          & 1        & 3 \\
Líneas G-code procesadas & 3\,262   & 205 \\
Factor de escala         & Fijo     & Configurable \\
Compensación Z           & Por zona Y & Interactiva \\
Manejo de errores        & Ninguno  & \texttt{limpiar\_robot()} \\
Parser regex             & Básico   & Robusto (negativos) \\
Cambio de herramienta    & No       & Sí (pausa + confirmar) \\
\bottomrule
\end{tabular}
\end{table}

% -------------------------------------------------------
\section{Discusión}

\subsection{Limitaciones del Sistema}

\begin{enumerate}
    \item \textbf{Velocidad de trazado:} A 100~mm/s y con comandos bloqueantes por punto, el sistema no escala bien para dibujos con decenas de miles de trayectorias. Una mejora futura sería utilizar la API de trayectorias (\texttt{set\_position} en modo FIFO o el comando nativo de G-code del SDK) para reducir la latencia entre puntos.

    \item \textbf{Planaridad asumida:} El sistema asume que el papel está en un plano fijo. Variaciones de más de 2~mm en la planaridad de la superficie producen trazos incompletos. Un sensor de fuerza en el efector permitiría implementar control de contacto.

    \item \textbf{G-code parcialmente soportado:} El intérprete actual ignora arcos (G02/G03) y solo procesa las coordenadas XY como movimientos lineales. Esto introduce un error geométrico en curvas, visible como una ligera poligonización de las partes circulares del dibujo.
\end{enumerate}

\subsection{Trabajo Futuro}

\begin{itemize}
    \item Implementar soporte completo para G02/G03 (interpolación circular) en el parser.
    \item Integrar el módulo \texttt{xarm.tools.gcode} del SDK para una ejecución de G-code nativa y más eficiente.
    \item Explorar la generación automática de archivos por capa desde Inkscape usando scripts de extensión en Python, eliminando el paso manual de exportación.
\end{itemize}

% -------------------------------------------------------
\section{Conclusiones}

Este proyecto demostró la viabilidad de integrar herramientas de diseño estándar (Inkscape, Gcodetools) con robots colaborativos industriales (Lite~6) mediante Python como lenguaje de orquestación. La evolución entre las dos versiones del controlador ---de \texttt{robot\_draw.py} a \texttt{colors-draw.py}--- ilustra un patrón de desarrollo iterativo donde cada iteración resuelve las limitaciones prácticas identificadas en la anterior: compensación de plano, soporte multicolor, manejo de errores y calibración interactiva.

La correcta parametrización del origen de coordenadas, el factor de escala y la orientación del efector fue tan determinante para el éxito del sistema como la propia cinemática del robot. Esto subraya que en sistemas Cobot aplicados a tareas de precisión, la robustez del software de control y la calidad del proceso de calibración son factores igualmente críticos a las especificaciones mecánicas del hardware.

% -------------------------------------------------------
\begin{thebibliography}{9}

\bibitem{murata_smr}
S.~Murata y H.~Kurokawa,
\textit{Self-Reconfigurable Robots: An Introduction}.
MIT Press, 2012.

\bibitem{iso15066}
ISO/TS~15066:2016,
``Robots and Robotic Devices --- Collaborative Robots,''
International Organization for Standardization, Geneva, Switzerland, 2016.

\bibitem{grieves_dt}
M.~Grieves y J.~Vickers,
``Digital Twin: Mitigating Unpredictable, Undesirable Emergent Behavior in Complex Systems,''
en \textit{Transdisciplinary Perspectives on Complex Systems},
Springer, 2017, pp.~85--113.

\bibitem{ufactory_lite6}
UFactory,
``Lite 6 User Manual v2.0,''
UFactory Co.~Ltd., Shenzhen, China, 2024.
[En línea]. Disponible: \url{https://www.ufactory.cc}

\bibitem{xarm_sdk}
UFactory,
``xArm Python SDK --- GitHub Repository,''
2024.
[En línea]. Disponible: \url{https://github.com/xArm-Developer/xArm-Python-SDK}

\bibitem{gcodetools}
N.~Russkikh,
``Gcodetools --- Inkscape G-code Extension,''
2023.
[En línea]. Disponible: \url{https://github.com/cnc-club/gcodetools}

\end{thebibliography}

\end{document}
